%% Đề 6 -- de-cuoi-hoc-ky-1-toan-9-thcs-nam-2024-2025-so-gddt-lang-son.pdf
\deheader{ĐỀ SỐ 6}%
  {Đề kiểm tra Học kỳ — Năm học 2024--2025}%
  {Sở GD\&ĐT Lạng Sơn | Chỉ phần Đại số}
\addcontentsline{toc}{section}{Đề 6 — Đề kiểm tra Học kỳ — Năm học 2024--2025}

\begin{de_cau}{Câu 1}{2,0 điểm}
\begin{enumerate}[label=\alph*)]
\item Tính: $A = \sqrt{50} - 3\sqrt{8} + \sqrt{18}$.
\item Rút gọn: $B = \left(\dfrac{2}{\sqrt{x}-2}+\dfrac{1}{\sqrt{x}+2}\right)\cdot\dfrac{\sqrt{x}-2}{\sqrt{x}+1}$ với $x\geq 0, x\neq 4$.
\end{enumerate}
\end{de_cau}

\begin{loi_giai_cau}
\textbf{a)} $A = 5\sqrt{2}-6\sqrt{2}+3\sqrt{2} = (5-6+3)\sqrt{2} = \boxed{2\sqrt{2}}$.

\textbf{b)}
\[B = \frac{2(\sqrt{x}+2)+(\sqrt{x}-2)}{(\sqrt{x}-2)(\sqrt{x}+2)}\cdot\frac{\sqrt{x}-2}{\sqrt{x}+1}
= \frac{3\sqrt{x}+2}{x-4}\cdot\frac{\sqrt{x}-2}{\sqrt{x}+1}
= \frac{3\sqrt{x}+2}{(\sqrt{x}+2)(\sqrt{x}+1)}.\]
\end{loi_giai_cau}

\begin{de_cau}{Câu 2}{2,0 điểm}
\begin{enumerate}[label=\alph*)]
\item Giải phương trình: $x^2 - 5x + 6 = 0$.
\item Giải hệ phương trình: 
\[ \begin{cases} x - 2y = 3 \\ 2x + y = 1 \end{cases} \]
\item Vẽ đồ thị hàm số $y = 2x^2$ (lập bảng giá trị với $x\in\{-2;-1;0;1;2\}$).
\end{enumerate}
\end{de_cau}

\begin{loi_giai_cau}
\textbf{a)} $x^2-5x+6=0\Rightarrow(x-2)(x-3)=0\Rightarrow x=2$ hoặc $x=3$.

\textbf{b)} Từ (1): $x=2y+3$. Thế vào (2): $2(2y+3)+y=1\Rightarrow 5y=-5\Rightarrow y=-1, x=1$.

Nghiệm $(x;y)=(1;-1)$.

\textbf{c)} Bảng giá trị:
\begin{center}
\begin{tabular}{c|ccccc}$x$&$-2$&$-1$&$0$&$1$&$2$\\\hline$y=2x^2$&$8$&$2$&$0$&$2$&$8$\end{tabular}
\end{center}
Đồ thị là Parabol có đỉnh $O(0;0)$, bề lõm quay lên.
\end{loi_giai_cau}

\begin{de_cau}{Câu 3}{1,5 điểm}
Một người đi xe máy từ A đến B với vận tốc 40 km/h. Khi về, người đó đi với vận tốc 50 km/h và về
sớm hơn lúc đi 15 phút. Tính độ dài quãng đường AB.
\end{de_cau}

\begin{loi_giai_cau}
Gọi độ dài $AB = d$ (km), $d>0$.

Thời gian đi: $\dfrac{d}{40}$ giờ. Thời gian về: $\dfrac{d}{50}$ giờ.

Về sớm hơn đi 15 phút $= \dfrac{1}{4}$ giờ:
\[\frac{d}{40} - \frac{d}{50} = \frac{1}{4}.\]
\[\frac{5d-4d}{200} = \frac{1}{4} \Rightarrow d = 50 \text{ km}.\]

Vậy quãng đường $AB = \boxed{50}$ km.
\end{loi_giai_cau}

\begin{de_cau}{Câu 4}{1,0 điểm}
Cho phương trình $x^2 - 2mx + m^2 - 1 = 0$ (tham số $m$). Tìm $m$ để phương trình có hai nghiệm
phân biệt $x_1, x_2$ thỏa mãn $x_1^2 + x_2^2 = 10$.
\end{de_cau}

\begin{loi_giai_cau}
$\Delta' = m^2 - (m^2-1) = 1 > 0$ với mọi $m$, phương trình luôn có hai nghiệm phân biệt.

Vi-ét: $x_1+x_2=2m$;\; $x_1x_2=m^2-1$.

\[x_1^2+x_2^2 = (x_1+x_2)^2-2x_1x_2 = 4m^2-2(m^2-1) = 2m^2+2 = 10.\]
\[m^2=4\Rightarrow\boxed{m=\pm 2}.\]
\end{loi_giai_cau}
